
\section{Security Definitions }
\label{app:secdef}
%\subsection{Unforgeability of \blklistToken}
%\vspace{-5pt}
%\label{app:defn:unforge}


\subsection{Weak Random Key Robustness}
\label{app:secdef:random_key}
We will use a symmetric key encryption scheme defined by the triple of algorithms $\encryptionSchemeDefn$. $\encryptKeyGen(1^\secParam)$ is a randomized key generation algorithm that on input a security parameter outputs a key \genericSymmKey from the key space \encryptKeySpace; \encrypt is a randomized encryption algorithm that on input a key \genericSymmKey and a message \plaintext outputs a ciphertext $\ciphertext \assign \encrypt[\genericSymmKey](\plaintext)$; \decrypt is a deterministic algorithm that on input a key \genericSymmKey and ciphertext \ciphertext outputs the plaintext \plaintext. Correctness of the encryption scheme implies $\decrypt[\genericSymmKey](\encrypt[\genericSymmKey](\plaintext)) = \plaintext$. 

The security property we will need is random key robustness which implies that a ciphertext encrypted under a randomly chosen key, cannot be decrypted under a different, randomly chosen key. It is known that an encryption scheme that satisfies IND-CPA\$ security and ciphertext integrity has weak random key robustness.


\subsubsection{IND-CPA\$ security}
\label{sec:defn:ind_cpa}
%
This property implies that a PPT. adversary cannot distinguish between a ciphertext produced by an encryption scheme \encryptionSchemeDefn and a 
random string of equal length. This is defined using the following experiment 



	\smallskip
	\begin{minipage}[t]{0.5\columnwidth}
			\begin{tabbing}
			*** \= \kill
			$\Experiment{\indcpaRandLabel[0]}{\encryptionSchemeDefn}(\indcpaAdversary)$ \\
			\> $\genericSymmKey \gets \encryptKeyGen()$ \\
			\> $\indcpaGuess \gets 
			\indcpaAdversary^{\cciOracle{\genericSymmKey}}()$ \\
			\> return \indcpaGuess
		\end{tabbing}
	\end{minipage}
	\hspace{50pt}
	\begin{minipage}[t]{0.5\columnwidth}
			\begin{tabbing}
			*** \= \kill
			$\Experiment{\indcpaRandLabel[1]}{\encryptionSchemeDefn}(\indcpaAdversary)$ \\
			\> $\genericSymmKey \gets \encryptKeyGen()$ \\
			\> $\indcpaGuess \gets 
			\indcpaAdversary^{\cciOracle(\cdot)}()$ \\
			\> return \indcpaGuess
		\end{tabbing}
	\end{minipage}
	


	
	\noindent
	where the maximum is taken over all adversaries \indcpaAdversary that run in
    \oracleQueries queries
	to \cciOracle{}.
    


Here, \cciOracle{\genericSymmKey} is an oracle to which the
key \genericSymmKey is a fixed input.  In each invocation to it, the
adversary \indcpaAdversary provides an input plaintext \plaintext and 
the oracle returns $\encrypt[\genericSymmKey](\plaintext)$. The oracle $\cciOracle$ 
returns random strings of length equal to a ciphertext.  




The adversary's
advantage is defined as

	\begin{eqnarray*}
		& \Advantage{\indcpaRandLabel}{\encryptionSchemeDefn}(\indcpaAdversary) = \\
		& \left|\prob{\Experiment{\indcpaRandLabel[1]}{\encryptionSchemeDefn}(\indcpaAdversary) = 1} - 
		\prob{\Experiment{\indcpaRandLabel[0]}{\encryptionSchemeDefn}(\indcpaAdversary) = 1} \right| \\
		& \Advantage{\indcpaRandLabel}{\encryptionSchemeDefn}(\timeBound, \oracleQueries) = 
		\max_{\indcpaAdversary} \Advantage{\indcpaRandLabel}{\encryptionSchemeDefn}(\indcpaAdversary)
	\end{eqnarray*}

\noindent
where the maximum is taken over all adversaries \indcpaAdversary running
in time at most \timeBound and making at most \oracleQueries queries. 

\subsubsection{Ciphertext integrity}
\label{sec:defn:cipher_integrity}
%
This property implies that a PPT. adversary cannot produce a new ciphertext after observing encryptions of polynomially-many ciphertexts under a randomly-chosen key. This is defined using the following experiment 

\begin{center}
	\begin{minipage}{0.5\columnwidth}
		\begin{tabbing}
			*** \= *** \= \kill
			$\Experiment{\cciLabel}{\encryptionSchemeDefn}(\indcpaAdversary)$ \\
			\> $\genericSymmKey \gets \encryptKeyGen()$ \\
			\> $\ciphertext' \gets 
			\indcpaAdversary^{\cciOracle{\genericSymmKey}}$ \\
            \> Let $\genericSetVar$ be the set of all ciphertexts returned by the oracle \cciOracle{\genericSymmKey} \\
            \> If $\ciphertext' \neq \mathsf{reject} ~\land \ciphertext' \notin \genericSetVar$\\
			\> \> then return 1\\
            \> \> else return 0
		\end{tabbing}
	\end{minipage}
\end{center}

Here, \cciOracle{\genericSymmKey} is an oracle to which the
key \genericSymmKey is a fixed input.  In each invocation to it, the
adversary \indcpaAdversary provides an input plaintext \plaintext and 
the oracle returns $\encrypt[\genericSymmKey](\plaintext)$. The adversary produces $\ciphertext'$
and wins the experiment if $\ciphertext'$ is a valid ciphertext and it has not been produced by \cciOracle{\genericSymmKey} 
in any previous query. 

The adversary's
advantage is defined as

\begin{align*}
	\Advantage{\cciLabel}{\encryptionSchemeDefn}(\indcpaAdversary)
	& = \prob{\Experiment{\cciLabel}{\encryptionSchemeDefn}(\indcpaAdversary) = 1} \\
	\Advantage{\cciLabel}{\encryptionSchemeDefn}(\timeBound, \oracleQueries) 
	&
	= \max_{\indcpaAdversary} \Advantage{\cciLabel}{\encryptionSchemeDefn}(\indcpaAdversary)
\end{align*}
where the maximum is taken over all adversaries \indcpaAdversary running
in time at most \timeBound and making at most \oracleQueries queries.

\subsection{Pseudorandom function}
\label{sec:defn:prf}
%



\newcommand{\allPerms}[1]{\ensuremath{\setNotation{P}({#1})}\xspace}
\newcommand{\allFuncs}[1]{\ensuremath{\setNotation{F}({#1})}\xspace}
\newcommand{\prfGuess}{\ensuremath{\valNotation{\hat{b}}}\xspace}
\newcommand{\prfAdversary}{\ensuremath{\algNotation{A}_{\mathsf{prf}}}\xspace}
\newcommand{\forwardOracleQueries}{\ensuremath{q}\xspace}
\NewDocumentCommand{\PRFExp}{o}{\ensuremath{\Experiment{\prfLabel\IfValueT{#1}{-#1}}{\dyPRF}(\prfAdversary)}\xspace}
\newcommand{\PRFAdv}{\ensuremath{\Advantage{\prfLabel}{\dyPRF}(\timeBound, \oracleQueries)}\xspace}
\NewDocumentCommand{\prfLabel}{o}{%
	{\scriptsize\textsf{prf\IfNoValueF{#1}{-{#1}}}}\xspace}


\begin{definition}[Pseudorandom function]
	\label{dfn:prp}
	Let $\genericFnFamily: \genericFnKeyspace \times \genericFnDomain
	\rightarrow \genericFnDomain$ be a function family,
	$\allFuncs{\genericFnDomain}$ be the set of all functions with domain
	and range \genericFnDomain, and \prfAdversary be an algorithm that
	takes an oracle of type $\genericFnDomain \rightarrow
	\genericFnDomain$ and returns a bit.  Let
	
	\smallskip
	\begin{minipage}[t]{0.4\columnwidth}
		\begin{tabbing}
			***\=***\=\kill
			$\Experiment{\prfLabel[0]}{\genericFnFamily}(\prfAdversary)$ \\
			\> $\genericFn \sample \allFuncs{\genericFnDomain}$ \\
			\> $\prfGuess \leftarrow \prfAdversary^{\genericFn(\cdot)}()$ \\
			\> return $\prfGuess$
		\end{tabbing}
	\end{minipage}
	\hspace{50pt}
	\begin{minipage}[t]{0.4\columnwidth}
		\begin{tabbing}
			***\=***\=\kill
			$\Experiment{\prfLabel[1]}{\genericFnFamily}(\prfAdversary)$ \\
			\> $\genericFnKey \sample \genericFnKeyspace$, $\genericFn \gets 
			\genericFnFamily[\genericFnKey]$ \\
			\> $\prfGuess \leftarrow \prfAdversary^{\genericFn(\cdot)}()$ \\
			\> return $\prfGuess$
		\end{tabbing}
	\end{minipage}
	
	\smallskip
	\noindent Then,
	
	\begin{eqnarray*}
		& \Advantage{\prfLabel}{\genericFnFamily}(\prfAdversary) = \\
		& \left|\prob{\Experiment{\prfLabel[1]}{\genericFnFamily}(\prfAdversary) = 1} - 
		\prob{\Experiment{\prfLabel[0]}{\genericFnFamily}(\prfAdversary) = 1} \right| \\
		& \Advantage{\prfLabel}{\genericFnFamily}(\timeBound, \oracleQueries)
		= \max_{\prfAdversary} \Advantage{\prfLabel}{\genericFnFamily}(\prfAdversary)
	\end{eqnarray*}
	
	\noindent
	where the maximum is taken over all adversaries \prfAdversary that run in
	time at most \timeBound and make at most \forwardOracleQueries queries
	to \genericFn.
\end{definition}

\iffalse 
\newcommand{\prgAdversary}{\ensuremath{\algNotation{A}_{\mathsf{prg}}}\xspace}
\NewDocumentCommand{\prgLabel}{o}{%
	{\scriptsize\textsf{prg\IfNoValueF{#1}{-{#1}}}}\xspace}
\subsection{Pseudorandom Generator}


\begin{definition}[Pseudorandom Generator]
	\label{dfn:prp}
	Let $\genericFnFamily: \genericFnKeyspace 
	\rightarrow \genericFnDomain$ be a function family,
	$\allFuncs{\genericFnDomain}$ be the set of all functions with domain
	\genericFnKeyspace, and \prgAdversary be an algorithm that
	takes an oracle of type $\genericFnDomain \rightarrow
	\genericFnDomain$ and returns a bit.  Let
	
	\smallskip
	\begin{minipage}[t]{0.4\columnwidth}
		\begin{tabbing}
			***\=***\=\kill
			$\Experiment{\prgLabel[0]}{\genericFnFamily}(\prgAdversary)$ \\
			\> $\genericFn \sample \allFuncs{\genericFnDomain}$ \\
			\> $\genericRand \sample \prgKeySpace$ \\
			\> $\prfGuess \leftarrow \prgAdversary(\genericFn(\genericRand))$ \\
			\> return $\prfGuess$
		\end{tabbing}
	\end{minipage}
	\hspace{50pt}
	\begin{minipage}[t]{0.4\columnwidth}
		\begin{tabbing}
			***\=***\=\kill
			$\Experiment{\prgLabel[1]}{\genericFnFamily}(\prgAdversary)$ \\\
			\> $\genericRand \sample \genericFnRange$\\
			\> $\prfGuess \leftarrow \prgAdversary(\genericRand)$ \\
			\> return $\prfGuess$
		\end{tabbing}
	\end{minipage}
	
	\smallskip
	\noindent Then,
	
	\begin{eqnarray*}
		& \Advantage{\prgLabel}{\genericFnFamily}(\prgAdversary) = \\
		& \left|\prob{\Experiment{\prgLabel[1]}{\genericFnFamily}(\prgAdversary) = 1} - 
		\prob{\Experiment{\prgLabel[0]}{\genericFnFamily}(\prgAdversary) = 1} \right| \\
		& \Advantage{\prgLabel}{\genericFnFamily}(\timeBound)
		= \max_{\prgAdversary} \Advantage{\prgLabel}{\genericFnFamily}(\prfAdversary)
	\end{eqnarray*}
	
	\noindent
	where the maximum is taken over all adversaries \prgAdversary that run in
	time at most \timeBound 
\end{definition}
\fi 